\section{Introduction}
\label{sec:intro}
\vspace{-2pt}
Recent advances in Large Language Model (LLM)-based Multi-Agent Systems \citep{li2024survey} have enhanced Retrieval-Augmented Generation (RAG) \citep{lewis2020retrieval, gao2023retrieval, jiang2023active} through  collaborating specialized agents \citep{nguyen2025ma, liu2025hm, qian2024scaling, chen2025improving}, extending  tool use \citep{zhuang2023toolchain, schick2023toolformer, jin2025search} and memory integration \citep{hu2025memory, gutierrez2025rag, yang2025improving}. These systems are promising for complex and long-context reasoning, where effective coordination between retrieval and reasoning is critical. 

Despite their promise, existing multi-agent RAG frameworks are fundamentally constrained by static or weakly adaptive orchestration. Most prior work relies on fixed or sequential pipelines \citep{gamage2024multi, liu2025hm, yao2022react, qian2024scaling, chen2025improving}, which fails to accommodate query-dependent variation in reasoning and retrieval complexity, resulting in insufficient evidence collection or unnecessary retrieval overhead. Besides, these systems are highly susceptible to error propagation \citep{shen2025understanding, pei2025scope}, where mistakes compound across multi-turn interactions and long-horizon agent coordination. Since these errors are inherently compositional, reliance on global supervision signals (\eg, final answer correctness) \citep{chen2025improving, song2025r1old, song2025r1} provides limited attribution to specific failure sources, hindering targeted improvements. Moreover, mainstream training-based optimization remains costly and inflexible, typically requiring large trajectory datasets \citep{zhu2024information, chen2025improving, chen2025mao, zhu2024atm, jin2025search}, which constrains scalability and limits continual adaptation in dynamic knowledge environments \citep{tang2025adapting}. As query distributions and underlying corpora evolve, fixed reasoning and coordination strategies \citep{jeong2024adaptive} lead to stagnant performance and brittle generalization to novel queries and retrieval challenges.

To address these limitations, we argue that, an effective multi-agent RAG system must (i) adapt its coordination strategies to query complexity, (ii) enable targeted and role-specific agent improvement, (iii) mitigate error propagation, and (iv) support continual adaptation without costly retraining. To this end, we propose \hero, a \underline{\textbf{H}}ierarchical \underline{\textbf{E}}volvution \underline{\textbf{RA}}G framework that achieves these goals by \textit{distilling execution experience as in-context knowledge} and \textit{reshaping agent behaviors by prompt evolution} without parameter updates. 

\hero is organized as a three-layer hierarchy that jointly evolves global orchestration strategies and local agent behaviors using accumulated experiential and reflective knowledge (Figure .\ref{fig:hera}). At the top, a centralized \textit{orchestrator} generates query-specific execution plans in a single step. This holistic orchestration \citep{ke2026mas} enables the system to reason globally about the coordination of retrieval and reasoning, producing plans tailored to the complexity and requirements of each query. At the bottom, \textit{specialized agents} execute subtasks grounded in interaction with others. \hero employs role-aware prompt optimization to drive targeted improvements for each agent’s unique responsibilities, avoiding generic updates and mitigating errors arising from misaligned multi-agent behaviors. Bridging these layers, the \textit{experience library} captures reflective insights from both successful and failed execution trajectories. These insights provide diagnostic information \citep{agrawal2025gepa}, enabling semantic credit assignment for global learning and guiding agent-level refinement. Across \hero, continuously evolving experience acts as a compass: it informs the orchestrator to improve high-level coordination strategies and drives progressive, role-specific enhancement of agent behavior, improving robustness and reducing cascading errors.


Our contributions are summarized as follow: (i) We introduce \hero, a hierarchical framework enabling dual-level evolution - experience-guided orchestration and role-aware prompt adaptation, shifting policy optimization from parameter updates to experience-driven contextual adaptation. (ii) We design topology-based metrics to quantitatively characterize the evolution of agent interaction topologies and track dynamics. (iii) Empirical results on six knowledge-intensive benchmarks demonstrate that \hero outperforms recent SOTA by an average of 38.69\% while maintaining token efficiency. 
